\documentclass{article}
\usepackage{../math_template}

\author{Jonathan Kasongo}
\title{European Girls Math Olympiad 2023 --- P1/3}

\begin{document}
\maketitle

\begin{problem}{European Girls Math Olympiad 2023 --- P1/3}
There are $n\geq 3$ positive real numbers $a_1, a_2, \ldots, a_n$. For
each $1\leq i \leq n$ we let $b_i = \frac{a_{i-1} + a_{i+1}}{a_i}$
(here we define $a_0$ to be $a_n$ and $a_{n+1}$ to be $a_1$). Assume that
for all $i$ and $j$ in the range 1 to $n$, we have $a_i \leq a_j$ if and
only if $b_i \leq b_j$. \vspace{0.2cm}

Prove that $a_1 = a_2 = \cdots = a_n$.
\end{problem}

\begin{solution}{Solution}
Suppose, for the sake of contradiction, that there exists
$1 \leq i < j \leq n$ such that $a_i < a_j$. Let $a_M$ be the largest
number in $(a_n)$. Let $a_m$ be the smallest number in $(a_n)$. Note that
$b_m$ is the smallest number in $(b_n)$ and that $b_M$ is the largest
number in $(b_n)$, because of the fact that
$a_i \leq a_j \iff b_i \leq b_j$. Now observe that

$$
b_M = \frac{a_{M-1} + a_{M+1}}{a_M} \leq 2, \ \ \ \ \ \
b_m = \frac{a_{m-1} + a_{m+1}}{a_m} \geq 2
$$

and so $b_m = b_M = 2$, which means that $b_1 = b_2 = \cdots = b_n = 2$.
Now since we have equality over all $b_i$ the statement
$b_i \leq b_j \iff a_i \leq a_j$ is true for all pairs of integers $(i,j)$.
That means that we have $a_i \leq a_j$ and $a_i \geq a_j$ and so $a_i = a_j$
again for all pairs of integers $(i,j)$. Finally this means that
$a_1 = a_2 = \cdots = a_n$ which is a contradiction. So indeed
$a_1 = a_2 = \cdots = a_n$. $\blacksquare$
\end{solution}


\end{document}
